\setcounter{section}{1}
\setcounter{subsection}{0}
\setcounter{theorem}{0}
\setcounter{figure}{1}
\subsection{Characterization via Prime Modulus} \label{sec:intro-char}

  We show, in \autoref{sec:characterization}, that every class $\PPA_q$ is built
out of the classes $\PPA_p$ for $p$ a prime. To formalize this result, we
recall the operator `\amp' defined by Buss and
Johnson \cite[\S6]{r13}. For any two syntactic complexity
classes $\M_0$, $\M_1$ with complete problems $S_0$, $S_1$, the class
$\M_0 \bamp \M_1$ is defined via its complete problem $S_0 \bamp S_1$ where, on
input $(x, b) \in \{0, 1\}^* \times \bit$, the goal is to find a solution for
$x$ interpreted as an instance of problem $S_b$. Namely, if $b = 0$ then the output has to be a solution of $S_0$ with
input $x$, and otherwise it has to be a solution of $S_1$ with input $x$.
Intuitively speaking, $\M_1 \bamp \M_2$ combines
the powers of both $\M_1$ and $\M_2$. Note that
$\M_1 \cup \M_2 \subseteq \M_1 \bamp \M_2$. We can now formally
express our characterization result (where $p|q$ is the set of primes $p$
dividing $q$).

\begin{theorem}\label{thm:prime-characterization}
$\PPA_q = \amp_{p|q}\, \PPA_p$.
\end{theorem}

\noindent A special case of \autoref{thm:prime-characterization} is that
$\PPA_{p^k}=\PPA_p$ for every prime power $p^k$. Showing the inclusion
$\PPA_{p^k}\subseteq \PPA_p$ is the crux of our proof. This part of the theorem can be viewed as a total search problem
analog of the counting class result of Beigel and Gill~\cite{r7}
stating that $\mathsf{Mod}_{p^k}\P = \mathsf{Mod}_{p}\P$; ``an unexpected
result'', they wrote at the time. Throughout this paper, we use $q$ to denote any integer $\ge 2$ and $p$ to denote a prime integer.

\paragraph{Concurrent independent work (original attribution)}
We note that a concurrent and independent work by Hollender~\cite{r25} also establishes the structural properties corresponding to the prime-modulus characterization and relative structural-power discussion.
